\documentclass[11pt]{amsart}
\usepackage[margin=1.1in]{geometry}
\usepackage{amsmath,amssymb,amsthm,mathtools}
\usepackage{booktabs}
\usepackage{hyperref}
\hypersetup{colorlinks=true,linkcolor=blue,citecolor=blue,urlcolor=blue}

\newtheorem{theorem}{Theorem}
\renewcommand{\thetheorem}{\Alph{theorem}}
\newtheorem{lemma}{Lemma}[section]
\newtheorem{proposition}[lemma]{Proposition}
\newtheorem{corollary}[lemma]{Corollary}
\theoremstyle{definition}
\newtheorem{definition}[lemma]{Definition}
\newtheorem{problem}[lemma]{Problem}
\theoremstyle{remark}
\newtheorem{remark}[lemma]{Remark}
\newtheorem{example}[lemma]{Example}

\newcommand{\Z}{\mathbb{Z}}
\newcommand{\Q}{\mathbb{Q}}
\newcommand{\F}{\mathbb{F}}
\newcommand{\PP}{\mathbb{P}}
\newcommand{\Gm}{\mathbb{G}_{\mathrm m}}
\newcommand{\R}{\mathcal{R}}
\DeclareMathOperator{\ord}{ord}
\DeclareMathOperator{\lcm}{lcm}

\title[Squares in $S$-unit shifts]{Exponent collapse and a question of Corvaja and Zannier\\ on squares in $S$-unit shifts}
\author{Keio}
\date{25 October 2026}

\begin{document}

\begin{abstract}
In a recent paper Corvaja and Zannier asked whether there are ``many'' integers $t_0$ such that $t_0+a2^m+b3^n$ is not a square for any integers $m,n$.
They remarked that congruence considerations might help, but that already for general $S$-units they did not see a possible path.
We prove that the integers $t_0$ for which some $t_0+a2^m+b3^n$ $(m,n\in\Z)$ is the square of a rational number form a set of natural density zero.
The same holds for sums of any number of terms from an arbitrary finitely generated subgroup of $\Q^\times$.
More generally, it holds for the integers represented by $F(x,z)$ with $x\in\Q$ and $z$ in a finitely generated subgroup of $(\Q^\times)^k$, whenever $F$ is a polynomial of even degree in $x$ whose leading coefficient is a monomial in $z$.
The proof adapts an argument of Bert\'ok and Hajdu for purely exponential equations.
For all primes $\ell$ with $\ell-1\mid n$ simultaneously, the group collapses to its finite quotient by $n$-th powers, whose order is polynomial in $n$; we call this ``exponent collapse''.
By a theorem of Prachar, the number of such primes can be made larger than any power of $\log n$.
The new ingredient is a value-set bound for the unknown $x\in\PP_1(\Q)$.
For $t_0+a2^m+b3^n$ with integers $a,b$ everything can be made explicit.
One squarefree modulus confines the represented $t_0$ to residue classes of total density less than $2.7\cdot10^{-22}$, uniformly in $a,b$.
For $a=b=1$ we also exhibit explicit arithmetic progressions in which no $t_0+2^m+3^n$ is a rational square, for example the class of $58$ modulo an explicit $37$-digit modulus.
\end{abstract}

\maketitle

\section{Introduction}

\subsection{The question}
Let $E/\Q$ be an elliptic curve, let $F\in\Q(x,z_1,z_2)$ and consider the values $F(x(P),2^m,3^n)$ with $P\in E(\Q)$ and $m,n\in\Z$.
Corvaja and Zannier \cite[Example~1.3]{CZ26} deduce from their Theorem~1.1 that infinitely many $t_0\in\Q$ are not of this form.
They then remark that replacing $E$ by $\PP_1$ ``would generally lead to problems which to our knowledge are difficult and widely open''.
For instance, the case $f(x,z_1,z_2)=x^2-az_1-bz_2$ leads to the question:

\begin{problem}[Corvaja--Zannier {\cite[after Example~1.3]{CZ26}}]\label{prob:CZ}
Are there ``many'' integers $t_0$ for which there are no integers $m,n$ such that $t_0+a2^m+b3^n$ is a square?
\end{problem}

They add that probabilistic considerations suggest that there are infinitely many such $t_0$, that ``possibly, congruence considerations can help here'', and that ``already replacing $2^m,3^n$ by more general $S$-units we do not see a possible path'' (they mention a conditional result of Bresciani \cite{Bre25}, which assumes conjectures of Bombieri--Lang type).

Here $m,n$ range over $\Z$, so $t_0+a2^m+b3^n\in\Z[1/6]$ and ``square'' means the square of a rational number; we answer the question in this form, which contains the variants with $m,n\ge0$ and integer squares.

\subsection{Results}
Throughout, density means natural density in $\Z$: a set $A\subseteq\Z$ has density zero if $\#\{t\in A:|t|\le X\}=o(X)$.
For $z=(z_1,\dots,z_k)$ and $\alpha\in\Z^k$ we write $z^\alpha=\prod_jz_j^{\alpha_j}$, and $\Q[z^{\pm1}]$ denotes the ring of Laurent polynomials in $z_1,\dots,z_k$.

\begin{definition}
For a subset $\Gamma\subseteq(\Q^\times)^k$ and $F\in\Q[z^{\pm1}][x]$ let
\[
\R(F,\Gamma)=\{\,t\in\Z:\ F(x,z)=t\ \text{for some }x\in\Q,\ z\in\Gamma\,\}.
\]
\end{definition}

\begin{theorem}\label{thm:A}
Let $\Gamma\subseteq(\Q^\times)^k$ be a finitely generated subgroup and let
\[
F(x,z)=c\,z^\alpha x^d+\sum_{i=0}^{d-1}A_i(z)\,x^i,\qquad c\in\Q^\times,\ \alpha\in\Z^k,\ A_i\in\Q[z^{\pm1}],
\]
with $d\ge2$ even.
Then $\R(F,\Gamma)$ has density zero.
More precisely, for every $\varepsilon>0$ there are a squarefree integer $M$ and a set $B\subseteq\Z/M\Z$ with $|B|\le\varepsilon M$ such that $t\bmod M\in B$ for every $t\in\R(F,\Gamma)$.
For $d=2$ the proof uses only Prachar's theorem (Lemma~\ref{lem:prachar}) and the number of squares in $\F_\ell$.
\end{theorem}

Taking $F=x^2-\sum_ia_iz_i$ we obtain the answer to Problem~\ref{prob:CZ}, for arbitrary $S$-units of $\Q$.

\begin{corollary}\label{cor:CZ}
Let $\Gamma_0\subseteq\Q^\times$ be finitely generated and $a_1,\dots,a_k\in\Q^\times$.
For all integers $t_0$ outside a set of density zero, $t_0+a_1u_1+\dots+a_ku_k$ is not the square of a rational number for any $u_1,\dots,u_k\in\Gamma_0$.
In particular, for all $t_0$ outside a set of density zero there are no $m,n\in\Z$ such that $t_0+a2^m+b3^n$ is a square in $\Q$.
\end{corollary}

The proof gives more than density.
For every $\varepsilon>0$ it gives a single squarefree modulus $M$, coprime to $6$ and to the denominators of $a,b$, with the following property.
For all $t_0$ outside residue classes of total density $\varepsilon$, the congruence $x^2\equiv t_0+a2^m+b3^n\pmod M$ has no solution with $x\in\Z$ and $m,n\in\Z$.
Hence, along a fixed sequence of moduli $M_1,M_2,\dots$, almost every $t_0$ is obstructed modulo some $M_j$.
This is a statement of Bert\'ok--Hajdu type (``the Hasse principle holds for almost all right-hand sides'') for this polynomial-exponential equation (Remark~\ref{rem:local}).
No single modulus obstructs all non-represented $t_0$, however.
The set $B$ of bad residues constructed in the proof is non-empty, since each $B_\xi$ has $\prod_\ell|V_{\xi,\ell}|\ge1$ elements.
Each class in $B$ has density $1/M>0$, while the represented integers have density zero.
So each class in $B$ also contains non-represented integers, and these are not obstructed modulo $M$.
A version over number fields would need Prachar-type results for primes of degree one, which we do not pursue.

For $t_0+a2^m+b3^n$ with $a,b\in\Z$ everything can be made explicit.

\begin{theorem}\label{thm:C}
For $a,b\in\Z$ let $\R_{a,b}$ be the set of $t_0\in\Z$ such that $t_0+a2^m+b3^n$ is a square in $\Q$ for some $m,n\in\Z$.
\begin{enumerate}
\item Let $E=1441440$ and let $L$ be the set of all prime factors of $\gcd(2^E-1,3^E-1)$.
These are $114$ primes, all larger than $3$, listed in Section~\ref{sec:explicit}.
Then for all $a,b\in\Z$ the set $\R_{a,b}$ is contained in a union of residue classes modulo $\prod_{\ell\in L}\ell$ of density less than $2.7\cdot10^{-22}$.
\item Let $M_{58}=5\cdot7\cdot11\cdot13\cdot19\cdot29\cdot31\cdot37\cdot41\cdot43\cdot61\cdot73\cdot127\cdot181\cdot211\cdot281\cdot337\cdot421\cdot631\cdot1009$, a $37$-digit number.
If $t\equiv58\pmod{M_{58}}$, then $t+2^m+3^n$ is not a square in $\Q$ for any $m,n\in\Z$.
Similar progressions through $t_0=-44$, $-38$, $-14$ and $100$ are given in Table~\ref{tab:certs}.
\item In the range $-50\le t_0\le100$ exactly five integers lie outside $\R_{1,1}$, namely $-44$, $-38$, $-14$, $58$ and $100$.
Every other $t_0$ in this range has a representation $t_0+2^m+3^n=x^2$ with $\max(|m|,|n|)\le9$.
\end{enumerate}
\end{theorem}

\subsection{The method}
Fix $n\ge1$ and consider all primes $\ell$ with $\ell-1\mid n$ (avoiding finitely many).
For every such $\ell$ and every $\gamma\in\Gamma$ one has $\gamma^n\equiv1\pmod\ell$ by Fermat's little theorem.
Hence, for all these primes simultaneously, the reduction of $\Gamma$ factors through the finite group $\Gamma/\Gamma^n$, of order at most $2^kn^{r}$ where $r=\operatorname{rank}\Gamma$.
For each of these $O(n^r)$ classes the polynomial $x\mapsto F(x,z)\bmod\ell$ has degree $d$ and misses a positive proportion of $\F_\ell$, so for each class the admissible $t\bmod\prod\ell$ have density at most $(1-\frac1{2d})^{\#\{\ell\}}$.
Prachar's theorem \cite{Pra55} provides $n$ for which the number of primes $\ell$ with $\ell-1\mid n$ exceeds $\exp(c\log n/(\log\log n)^2)$, which beats the polynomial number of classes.
We call this \emph{exponent collapse}.

The device itself is classical, and ``exponent collapse'' is only our name for it.
Moduli built from primes with $\ell-1\mid n$, equivalently moduli with small Carmichael function, are a standard tool for exponential equations.
They appear in work on Skolem's exponential local--global principle \cite{Sko37,BBL13}.
They are also the key to the theorem of Bert\'ok and Hajdu \cite{BH16,BH18} that the corresponding Hasse principle holds for almost all right-hand sides of a purely exponential equation; that theorem uses the result of Erd\H{o}s, Pomerance and Schmutz \cite{EPS91}.
Theorem~\ref{thm:A} is an adaptation of the Bert\'ok--Hajdu argument.
The new points here are the following.
\begin{itemize}
\item The unknown $x\in\PP_1(\Q)$ enters through the value set of a polynomial over $\F_\ell$. This is what allows the $\PP_1$-variable of Problem~\ref{prob:CZ}; for $d>2$ the value set is controlled by Carlitz's conjecture, proved by Fried, Guralnick and Saxl \cite{FGS93}, and by the value-set bound of Wan \cite{Wan93,WSC93}.
\item The argument is applied to the problem of Corvaja and Zannier.
\item Everything is made explicit for $t_0+2^m+3^n$.
\end{itemize}
For a fixed $z$ the problem is easy: if $F(x,z)\in\Z$ then the denominator of $x$ is bounded in terms of $z$, so $F(\Q,z)\cap\Z$ has $O_z(X^{1/d})$ elements of absolute value at most $X$.
But $\R(F,\Gamma)$ is a union over infinitely many $z$, and the bound is not uniform in $z$.
Theorem~\ref{thm:A} says that such a union, parametrised by a finitely generated group, still has density zero.

\subsection*{Organisation}
Section~\ref{sec:lemmas} collects the lemmas and Section~\ref{sec:proofA} proves Theorem~\ref{thm:A} and Corollary~\ref{cor:CZ} in an explicit form (Proposition~\ref{prop:explicit}).
Section~\ref{sec:explicit} proves Theorem~\ref{thm:C}.
Section~\ref{sec:remarks} discusses the hypotheses and open questions.
All computations are exact.
Theorem~\ref{thm:C}(1) is an exact rational computation.
The certificates of Theorem~\ref{thm:C}(2) were found by a vectorised search and re-verified by a separate pure-Python program.
Theorem~\ref{thm:C}(3) combines these certificates with an exhaustive search for representations with $\max(|m|,|n|)\le12$.
The code and data are in the supplementary material.

\section{Lemmas}\label{sec:lemmas}

For $n\ge1$ let $\omega^*(n)$ denote the number of primes $\ell$ with $(\ell-1)\mid n$.

\begin{lemma}[Prachar {\cite{Pra55}}]\label{lem:prachar}
There is an absolute constant $c_1>0$ such that
\[
\omega^*(n)>\exp\bigl(c_1\log n/(\log\log n)^2\bigr)\qquad\text{for infinitely many }n.
\]
\end{lemma}

An improvement to $\exp(c_2\log n/\log\log n)$ with $c_2=0.6736\log2$ is proved by Fan and Pollack \cite{FP25}.
We only need $\omega^*(n)/\log n\to\infty$ along some sequence of $n$.

\begin{lemma}[Exponent collapse]\label{lem:collapse}
Let $\Gamma\subseteq(\Q^\times)^k$ be generated by $\gamma_1,\dots,\gamma_s$ and let $S_0$ be the set of primes dividing the numerator or denominator of some coordinate of some $\gamma_i$.
Let $n\ge1$ and let $\ell\notin S_0$ be a prime with $\gamma_i^n\equiv(1,\dots,1)\pmod\ell$ for all $i$; this holds in particular if $(\ell-1)\mid n$.
Then every coordinate of every $\gamma\in\Gamma$ is an $\ell$-adic unit, and the reduction map $\rho_\ell:\Gamma\to(\F_\ell^\times)^k$ is trivial on $\Gamma^n=\{\gamma^n:\gamma\in\Gamma\}$.
Moreover $|\Gamma/\Gamma^n|\le2^kn^r$, where $r$ is the rank of $\Gamma$.
\end{lemma}

\begin{proof}
The coordinates of the generators are $\ell$-adic units, hence so are those of every element of $\Gamma$, and $\rho_\ell$ is a homomorphism.
Since $\rho_\ell(\gamma_i)^n=1$ for every $i$, $\rho_\ell$ kills $\Gamma^n$.
If $(\ell-1)\mid n$, then $u^n\equiv1\pmod\ell$ for every $\ell$-adic unit $u$ by Fermat's little theorem.
Finally, $\Gamma\cong T\times\Z^r$ with $T$ finite.
Since the torsion subgroup of $(\Q^\times)^k$ is $\{\pm1\}^k$, we have $|T|\le2^k$.
Since $\Gamma^n\cong T^n\times(n\Z)^r$, we get $|\Gamma/\Gamma^n|=|T/T^n|\,n^r\le2^kn^r$.
\end{proof}

\begin{lemma}[Integrality]\label{lem:integral}
Let $\Gamma$ and $S_0$ be as in Lemma~\ref{lem:collapse}.
Let $F$ be as in Theorem~\ref{thm:A}, except that $d\ge1$ is arbitrary.
Let $S_1\supseteq S_0$ also contain the primes dividing the numerator or the denominator of $c$, and the primes dividing the denominator of a coefficient of some $A_i$.
If $\ell\notin S_1$ and $F(x,z)=t$ with $x\in\Q$, $z\in\Gamma$ and $t\in\Z$, then $x$ is $\ell$-integral and $t\equiv\overline F(\bar x,\bar z)\pmod\ell$.
Here bars denote reduction modulo $\ell$, and $\overline F(x,\bar z)\in\F_\ell[x]$ has degree exactly $d$.
\end{lemma}

\begin{proof}
All $z_j$ are $\ell$-adic units, so $c\,z^\alpha$ is an $\ell$-adic unit and every $A_i(z)$ is $\ell$-integral.
In particular the leading coefficient $\bar c\,\bar z^\alpha$ of $\overline F(x,\bar z)$ is non-zero.
Suppose $v_\ell(x)=-s<0$.
Then $v_\ell(c\,z^\alpha x^d)=-ds$, while $v_\ell(A_i(z)x^i)\ge-is>-ds$ for $i<d$.
Hence $v_\ell(F(x,z))=-ds<0$, contradicting $F(x,z)=t\in\Z$.
So $x$ is $\ell$-integral, and reducing the identity $F(x,z)=t$ modulo $\ell$ gives the congruence.
\end{proof}

\begin{lemma}[Value sets]\label{lem:values}
Let $\ell$ be an odd prime, let $d\ge2$ and let $G\in\F_\ell[x]$ have degree $d$.
\begin{enumerate}
\item If $d=2$, then $|G(\F_\ell)|=(\ell+1)/2$.
\item If $d$ is even, there is a constant $C(d)$ such that $|G(\F_\ell)|\le\ell-(\ell-1)/d$ whenever $\ell>C(d)$.
\end{enumerate}
In both cases $|G(\F_\ell)|\le(1-\tfrac1{2d})\,\ell$.
\end{lemma}

\begin{proof}
(1) Write $G=u(x+v)^2+w$ with $u\ne0$.
The set $\{y^2:y\in\F_\ell\}$ has $(\ell+1)/2$ elements, and $y\mapsto uy+w$ is a bijection of $\F_\ell$.

(2) Carlitz's conjecture, proved by Fried, Guralnick and Saxl \cite{FGS93}, says that for every even $d$ there is $C(d)$ such that no polynomial of degree $d$ permutes $\F_q$ when $q$ is odd and $q>C(d)$.
(More generally: Lenstra proved that an exceptional polynomial of degree $d$ over $\F_q$ satisfies $\gcd(d,q-1)=1$, and Cohen and Fried gave an elementary version of his proof \cite{CF95}. A permutation polynomial of degree $d$, with $d$ prime to the characteristic, over a sufficiently large $\F_q$ is exceptional \cite{vzG91}. Together these give the Carlitz--Wan conjecture: no permutation polynomials of degree $d$ exist when $\gcd(d,q-1)>1$ and $q$ is large.)
So for $\ell>C(d)$ the polynomial $G$ is not a permutation polynomial.
By a theorem of Wan \cite{Wan93} (see also \cite[Theorem~3.1]{WSC93}), a polynomial of degree $d$ over $\F_q$ that is not a permutation polynomial takes at most $q-\lceil(q-1)/d\rceil$ values.

For the last assertion: in case (1), $(\ell+1)/(2\ell)\le2/3\le1-\frac1{2d}$ when $d=2$.
In case (2), $1-(\ell-1)/(d\ell)\le1-\frac1{2d}$ because $(\ell-1)/\ell\ge\frac12$.
\end{proof}

\section{Proof of Theorem~\ref{thm:A}}\label{sec:proofA}

We prove an explicit sieve inequality from which Theorem~\ref{thm:A} follows by Lemma~\ref{lem:prachar}.

\begin{proposition}\label{prop:explicit}
Let $\Gamma$, $F$, $S_0$ and $S_1$ be as in Theorem~\ref{thm:A} and Lemmas~\ref{lem:collapse} and~\ref{lem:integral}.
Put $S=S_1\cup\{2\}$ if $d=2$, and $S=S_1\cup\{2\}\cup\{\text{primes }\le C(d)\}$ if $d\ge4$, with $C(d)$ as in Lemma~\ref{lem:values}.
Let $n\ge1$, and let $L$ be a finite set of primes $\ell\notin S$ such that $\gamma^n\equiv1\pmod\ell$ for all $\gamma\in\Gamma$.
Put $M=\prod_{\ell\in L}\ell$, and let $\theta_\ell=(\ell+1)/(2\ell)$ if $d=2$ and $\theta_\ell=1-(\ell-1)/(d\ell)$ if $d\ge4$.
Then there is a set $B\subseteq\Z/M\Z$ such that $t\bmod M\in B$ for all $t\in\R(F,\Gamma)$, and
\[
\frac{|B|}{M}\ \le\ |\Gamma/\Gamma^n|\cdot\prod_{\ell\in L}\theta_\ell\ \le\ 2^kn^r\Bigl(1-\frac1{2d}\Bigr)^{|L|}.
\]
\end{proposition}

\begin{proof}
Let $\xi\in\Gamma/\Gamma^n$ and $\ell\in L$.
By Lemma~\ref{lem:collapse}, $\rho_\ell(z)$ is the same for all $z\in\xi$.
Hence $G_{\xi,\ell}(x):=\overline F(x,\rho_\ell(z))\in\F_\ell[x]$ is well defined.
It has degree $d$ by Lemma~\ref{lem:integral}.
Let $V_{\xi,\ell}=G_{\xi,\ell}(\F_\ell)$.
By the Chinese remainder theorem the set
\[
B_\xi=\{\tau\in\Z/M\Z:\ \tau\bmod\ell\in V_{\xi,\ell}\text{ for all }\ell\in L\}
\]
has $\prod_{\ell}|V_{\xi,\ell}|\le M\prod_\ell\theta_\ell$ elements.
Put $B=\bigcup_\xi B_\xi$.
If $t\in\R(F,\Gamma)$, say $F(x,z)=t$, let $\xi=z\Gamma^n$.
By Lemma~\ref{lem:integral}, $t\bmod\ell\in V_{\xi,\ell}$ for every $\ell\in L$, i.e. $t\bmod M\in B_\xi$.
The bound $|B|\le|\Gamma/\Gamma^n|\max_\xi|B_\xi|$ and Lemmas~\ref{lem:collapse} and~\ref{lem:values} give the inequality.
Here Lemma~\ref{lem:values} gives $|V_{\xi,\ell}|\le\theta_\ell\,\ell$, because $\ell\notin S$ is odd and, if $d\ge4$, larger than $C(d)$.
\end{proof}

\begin{proof}[Proof of Theorem~\ref{thm:A}]
Let $\varepsilon>0$.
For $n\ge1$ let $L_n=\{\ell\notin S:(\ell-1)\mid n\}$.
By Lemma~\ref{lem:collapse}, $L_n$ is admissible in Proposition~\ref{prop:explicit}, and $|L_n|\ge\omega^*(n)-|S|$.
By Lemma~\ref{lem:prachar} there are arbitrarily large $n$ with $\omega^*(n)>\exp(c_1\log n/(\log\log n)^2)$.
For these,
\[
2^kn^r\Bigl(1-\frac1{2d}\Bigr)^{\omega^*(n)-|S|}\longrightarrow0 ,
\]
since $\omega^*(n)/\log n\to\infty$.
Choose such an $n$ with this quantity at most $\varepsilon$ and let $M$, $B$ be as in Proposition~\ref{prop:explicit} for $L=L_n$.
Then $M$ is squarefree, $|B|\le\varepsilon M$, and $\R(F,\Gamma)\bmod M\subseteq B$.
Consequently $\#\{t\in\R(F,\Gamma):|t|\le X\}\le|B|\,\bigl((2X+1)/M+1\bigr)\le\varepsilon(2X+1)+M$ for all $X\ge1$.
So the upper density of $\R(F,\Gamma)$ is at most $\varepsilon$.
As $\varepsilon$ was arbitrary, $\R(F,\Gamma)$ has density zero.
\end{proof}

\begin{proof}[Proof of Corollary~\ref{cor:CZ}]
Apply Theorem~\ref{thm:A} with $d=2$, $F(x,z)=x^2-\sum_{i=1}^ka_iz_i$ and $\Gamma=\Gamma_0^k\subseteq(\Q^\times)^k$.
A relation $t_0+\sum_ia_iu_i=x^2$ with $x\in\Q$ and $u\in\Gamma_0^k$ means exactly $F(x,u)=t_0$.
For the second statement take $\Gamma_0=\langle2,3\rangle$, $k=2$, $(a_1,a_2)=(a,b)$, and $u_1=2^m$, $u_2=3^n$.
\end{proof}

\begin{remark}[Local obstruction]\label{rem:local}
The proof shows that every $t\in\R(F,\Gamma)$ satisfies a congruence condition modulo $M$.
For $t$ outside the set $B$ of residues modulo $M$, the congruence $F(x,z)\equiv t\pmod M$ has no solution with $x\in\Z/M\Z$ and $z\in\Gamma$.
This proves the statement after Corollary~\ref{cor:CZ}.
For $F=x^2-az_1-bz_2$ the set $S$ consists of $2$, $3$ and the primes dividing the denominator of $a$ or of $b$, so $M$ is coprime to all of them; for $a,b\in\Z$ only $2$ and $3$ are excluded.
\end{remark}

\section{Explicit computations for \texorpdfstring{$t_0+a2^m+b3^n$}{t0+a2m+b3n}}\label{sec:explicit}

Here $k=2$, $\Gamma=\langle2\rangle\times\langle3\rangle\subseteq(\Q^\times)^2$ and $F=x^2-az_1-bz_2$.
Then $S_0=\{2,3\}$ and $\Gamma/\Gamma^n\cong(\Z/n\Z)^2$, so $|\Gamma/\Gamma^n|=n^2$.
For $a,b\in\Z$ the set $S$ of Proposition~\ref{prop:explicit} is $\{2,3\}$.
So we may use every prime $\ell>3$ with $2^n\equiv3^n\equiv1\pmod\ell$, i.e.\ every prime factor of $\gcd(2^n-1,3^n-1)$; note that this $\gcd$ is prime to $6$.
Moreover $\theta_\ell=(\ell+1)/(2\ell)$ for every class and every $a,b$.
This is because $x\mapsto x^2-a2^\mu-b3^\nu$ takes exactly $(\ell+1)/2$ values modulo $\ell$, even when $\ell\mid ab$.

\begin{proof}[Proof of Theorem~\ref{thm:C}(1)]
Let $n=E=1441440=2^5\cdot3^2\cdot5\cdot7\cdot11\cdot13$.
The number $g=\gcd(2^E-1,3^E-1)$ has $379$ decimal digits.
Trial division by the primes below $3\cdot10^6$ factors it completely.
Its prime factors are the $114$ primes
\begin{quote}\small
5, 7, 11, 13, 17, 19, 23, 29, 31, 37, 41, 43, 53, 61, 67, 71, 73, 79, 89, 97, 113, 127, 131, 157, 181, 193, 199, 211, 241, 281, 313, 331, 337, 353, 397, 421, 463, 521, 547, 577, 617, 631, 661, 673, 859, 881, 911, 937, 991, 1009, 1093, 1171, 1249, 1321, 1873, 2003, 2017, 2081, 2113, 2311, 2341, 2521, 2731, 2861, 2971, 3121, 3169, 3361, 3433, 3697, 4621, 5281, 6007, 6337, 6553, 7393, 7489, 8009, 8191, 8581, 8737, 9241, 13441, 13729, 14561, 16381, 18481, 20021, 20161, 20593, 21841, 24571, 25741, 26209, 27457, 36037, 37441, 48049, 51481, 55441, 65521, 72073, 96097, 110881, 120121, 131041, 180181, 192193, 270271, 332641, 411841, 576577, 960961, 1921921.
\end{quote}
Among them, $96$ satisfy $(\ell-1)\mid E$.
The other $18$ (from $193$ to $1921921$) satisfy $(\ell-1)\nmid E$, but the orders of $2$ and $3$ modulo $\ell$ still divide $E$.
The program checks $2^E\equiv3^E\equiv1\pmod\ell$ for each $\ell$ and evaluates the rational number $E^2\prod_{\ell\in L}\frac{\ell+1}{2\ell}$ exactly.
The result is $2.6839\ldots\cdot10^{-22}<2.7\cdot10^{-22}$.
Since $S=\{2,3\}$ and $\theta_\ell=(\ell+1)/(2\ell)$ for all integers $a,b$, Proposition~\ref{prop:explicit} completes the proof.
For rational $a,b$ one removes from $L$ the primes dividing their denominators.
\end{proof}

\begin{proof}[Proof of Theorem~\ref{thm:C}(2)]
Let $P$ be a finite set of primes $\ell>3$ and let
\[
E_P=\lcm\{\ord_\ell(2),\ord_\ell(3):\ell\in P\}.
\]
Call $(t_0,P)$ a \emph{certificate} if for every $(\mu,\nu)\in(\Z/E_P\Z)^2$ there is an $\ell\in P$ such that $t_0+2^\mu+3^\nu$ is a non-zero quadratic non-residue modulo $\ell$.
Let $(t_0,P)$ be a certificate, let $t\equiv t_0\pmod{\prod_{\ell\in P}\ell}$, and let $m,n\in\Z$.
Put $\mu=m\bmod E_P$ and $\nu=n\bmod E_P$.
For the prime $\ell\in P$ attached to $(\mu,\nu)$ we have $t+2^m+3^n\equiv t_0+2^\mu+3^\nu\pmod\ell$.
This is a non-square modulo $\ell$.
If $t+2^m+3^n$ were equal to $x^2$ with $x\in\Q$, then $x$ would be $\ell$-integral, since $t+2^m+3^n$ is, and $t+2^m+3^n$ would be a square modulo $\ell$.
That is a contradiction.
Table~\ref{tab:certs} lists five certificates.
Each was found by a vectorised search over $(\Z/E\Z)^2$ with $E\in\{2520,5040\}$ and checked again by an independent pure-Python program, which runs through all $E_P^2$ classes; the data are in the supplementary material.
\end{proof}

\begin{table}[ht]
\centering
\caption{Certificates $(t_0,P)$ for $t_0+2^m+3^n$. For every $t\equiv t_0\pmod{\prod_{\ell\in P}\ell}$ and all $m,n\in\Z$, $t+2^m+3^n$ is not a square in $\Q$.}\label{tab:certs}
\small
\begin{tabular}{@{}rrrp{8.6cm}@{}}
\toprule
$t_0$ & $|P|$ & $E_P$ & $P$\\
\midrule
$58$ & 20 & 2520 & 5, 7, 11, 13, 19, 29, 31, 37, 41, 43, 61, 73, 127, 181, 211, 281, 337, 421, 631, 1009\\
$-44$ & 26 & 5040 & 5, 7, 11, 13, 17, 19, 29, 31, 37, 41, 43, 61, 71, 73, 97, 113, 127, 181, 211, 241, 281, 337, 421, 577, 631, 673\\
$-38$ & 25 & 5040 & same as for $-44$ without 211\\
$-14$ & 26 & 5040 & same as for $-44$\\
$100$ & 27 & 5040 & same as for $-44$, and 3361\\
\bottomrule
\end{tabular}
\end{table}

\begin{proof}[Proof of Theorem~\ref{thm:C}(3)]
The five values are excluded by the certificates of Table~\ref{tab:certs}.
For each of the other $146$ integers in $[-50,100]$ a search over $\max(|m|,|n|)\le12$ finds an explicit representation, with exact rational arithmetic.
The largest exponent needed is $9$, for example $37+2^9+3^3=24^2$.
\end{proof}

\begin{remark}[Small versus large $t_0$]
By Theorem~\ref{thm:C}(3), small $t_0$ are mostly represented: for instance $-1+2+3=2^2$, and $t_0=y^2-2$ with $m=n=0$.
This is not in conflict with Theorem~\ref{thm:C}(1), which concerns large $t_0$.
As a sanity check, $97$ of $100$ integers drawn at random from $[10^{15},10^{16})$ have certificates with $E=5040$.
For comparison, the bound of Proposition~\ref{prop:explicit} for $E=5040$ with the $31$ prime factors of $\gcd(2^{5040}-1,3^{5040}-1)$ is $0.0270$.
\end{remark}

\section{Remarks on the hypotheses, and open questions}\label{sec:remarks}

\begin{remark}[The leading coefficient]
Some hypothesis on the leading coefficient is necessary.
For $F=(z-1)x^2+x$ and $\Gamma=\langle2\rangle$ we have $F(x,1)=x$, so $\R(F,\Gamma)=\Z$.
The monomial hypothesis guarantees that $x\mapsto F(x,z)\bmod\ell$ has degree $d$ for every $z\in\Gamma$ and every $\ell\notin S$.
If the leading coefficient $A_d(z)$ is only assumed non-zero on $\Gamma$, a class $\xi\in\Gamma/\Gamma^n$ may be degenerate (that is, $A_d\equiv0$) modulo many $\ell$ at once.
Controlling such classes looks like a Skolem-type problem.
\end{remark}

\begin{remark}[One $\PP_1$-variable]
The method needs a single variable $x$.
For $F=x^2+y^2-z$ every residue modulo $\ell$ is a sum of two squares, so no prime gives information.
\end{remark}

\begin{remark}[Odd degree]
For $F=x^3-z$ and $\Gamma=\langle2\rangle$ the primes $\ell\equiv2\pmod3$ give no information, since cubing permutes $\F_\ell$.
Theorem~\ref{thm:A} extends to odd $d$ once one knows that, for a prime $q\mid d$ and every $A$, there are $n$ with more than $A\log n$ primes $\ell\equiv1\pmod q$ satisfying $(\ell-1)\mid n$.
This is a version of Prachar's theorem in arithmetic progressions.
The constructions of Carmichael numbers in arithmetic progressions (see e.g.\ \cite{Wri13}) produce many primes in prescribed residue classes with $\ell-1$ dividing a common number, and may well supply it.
We have not checked the details and do not claim the odd case.
\end{remark}

\begin{remark}[What cannot be improved]
The set $\R(F,\Gamma)$ is in general infinite.
For instance, if $a_1,\dots,a_k\in\Z$, the set in Corollary~\ref{cor:CZ} contains $y^2-\sum_ia_i$ for every $y\in\Z$ (take all $u_i=1$).
So ``density zero'' cannot be replaced by ``finite''.
The proof gives a rate of decay that depends on the quantitative form of Prachar's theorem, but no power saving.
\end{remark}

\begin{problem}
Let $F\in\Q[z^{\pm1}][x]$ have degree $d\ge2$ in $x$, and suppose that $F(x,z)$ has degree $d$ in $x$ for every $z\in\Gamma$.
Does $\R(F,\Gamma)$ have density zero?
What if $F$ is only assumed to be a rational function of degree at least $2$ in $x$?
\end{problem}

\begin{thebibliography}{99}
\bibitem{BBL13} B. Bartolome, Y. Bilu, F. Luca, On the exponential local-global principle, \emph{Acta Arith.} \textbf{159} (2013), 101--111.
\bibitem{BH16} Cs. Bert\'ok, L. Hajdu, A Hasse-type principle for exponential diophantine equations and its applications, \emph{Math. Comp.} \textbf{85} (2016), 849--860; arXiv:1407.6499.
\bibitem{BH18} Cs. Bert\'ok, L. Hajdu, A Hasse-type principle for exponential Diophantine equations over number fields and its applications, \emph{Monatsh. Math.} \textbf{187} (2018), 425--436.
\bibitem{Bre25} G. Bresciani, A higher dimensional Hilbert irreducibility theorem, \emph{Amer. J. Math.} (2025), doi:10.1353/ajm.2025.a961346; arXiv:2101.01090.
\bibitem{CF95} S. D. Cohen, M. D. Fried, Lenstra's proof of the Carlitz--Wan conjecture on exceptional polynomials: an elementary version, \emph{Finite Fields Appl.} \textbf{1} (1995), 372--375.
\bibitem{CZ26} P. Corvaja, U. Zannier, Rational and integral values of rational functions at rational points (with an appendix by D. Masser), arXiv:2608.28255 (2026).
\bibitem{EPS91} P. Erd\H{o}s, C. Pomerance, E. Schmutz, Carmichael's lambda function, \emph{Acta Arith.} \textbf{58} (1991), 363--385.
\bibitem{FGS93} M. D. Fried, R. Guralnick, J. Saxl, Schur covers and Carlitz's conjecture, \emph{Israel J. Math.} \textbf{82} (1993), 157--225.
\bibitem{FP25} K. Fan, P. Pollack, The maximal order of the shifted-prime divisor function, arXiv:2510.14167 (2025).
\bibitem{Pra55} K. Prachar, \"Uber die Anzahl der Teiler einer nat\"urlichen Zahl, welche die Form $p-1$ haben, \emph{Monatsh. Math.} \textbf{59} (1955), 91--97.
\bibitem{Sko37} Th. Skolem, Anwendung exponentieller Kongruenzen zum Beweis der Unl\"osbarkeit gewisser diophantischer Gleichungen, \emph{Avh. Norske Vid. Akad. Oslo I} (1937), no.~12.
\bibitem{vzG91} J. von zur Gathen, Values of polynomials over finite fields, \emph{Bull. Austral. Math. Soc.} \textbf{43} (1991), 141--146.
\bibitem{Wan93} D. Wan, A $p$-adic lifting lemma and its applications to permutation polynomials, in: \emph{Finite Fields, Coding Theory, and Advances in Communications and Computing}, Lecture Notes in Pure and Appl. Math. \textbf{141}, Dekker, New York, 1993, 209--216.
\bibitem{WSC93} D. Wan, P. J.-S. Shiue, C. S. Chen, Value sets of polynomials over finite fields, \emph{Proc. Amer. Math. Soc.} \textbf{119} (1993), 711--717.
\bibitem{Wri13} T. Wright, Infinitely many Carmichael numbers in arithmetic progressions, \emph{Bull. London Math. Soc.} \textbf{45} (2013), 943--952; arXiv:1212.5850.
\end{thebibliography}

\end{document}
